%% [AI DRAFT] Borrador generado para discusión (Modo C, etapa 3).
%% Citas verificadas contra el PDF. Sin marcadores [VERIFY] pendientes.

\section{CONCLUSIONES Y TRABAJO FUTURO}

Este trabajo ha presentado el diseño de un sistema que estima el riesgo de estafa en el punto
de decisión de compra, a partir de características que el comprador ya tiene delante y sin acceso
privilegiado a la plataforma. Tres conclusiones se desprenden del diseño.

La primera es que la evaluación de riesgo puede desplazarse desde el juicio del comprador
hacia el sistema sin requerir la cooperación de la plataforma: las ocho características de la
Tabla~\ref{tab:senales} se derivan del contenido que el navegador ya recibió, lo que distingue
la propuesta de la línea dominante en detección de fraude, que alcanza desempeños elevados
pero requiere telemetría interna \cite{ref28}.

La segunda es que la regla de abstención explícita no es un detalle de implementación sino una
condición de diseño. Un instrumento que presentara como seguro aquello que solo desconoce
trasladaría al comprador un riesgo que no ha evaluado, y penalizaría a comercios legítimos que
simplemente no publicitan sus garantías.

La tercera es que intervenir sobre la información del comprador tiene fundamento teórico y no
solo utilidad práctica: la cantidad de reseñas falsas en equilibrio decrece con la fracción de
consumidores capaces de interpretar correctamente los indicadores de una ficha \cite{ref29},
de modo que el efecto de la herramienta no se agota en quien la instala.

La limitación principal ha sido declarada a lo largo del artículo: la arquitectura
no ha sido implementada. El único resultado experimental disponible es de naturaleza
diagnóstica y apunta contra los propios datos: cinco algoritmos clasifican sin error la
partición de prueba, lo que revela que el conjunto es separable por construcción y que las
métricas obtenidas sobre él no estiman la calidad de la estimación de riesgo. Lo presentado es un diseño y un protocolo para
someterlo a prueba, no un sistema evaluado.

El trabajo futuro tiene un orden obligado, porque cada paso condiciona la interpretación del
siguiente. El primero es medir qué características resultan legibles en fichas reales. Importa
menos el porcentaje global que su reparto entre plataformas grandes y vendedores
independientes: si la cobertura se concentrara en las primeras, que son las que ya ofrecen
garantías al comprador, el sistema fallaría precisamente donde este se encuentra más expuesto.
El segundo es construir un conjunto de datos más sólido, cuyas etiquetas provengan de una
fuente independiente del vector de entrada y cuyo tamaño y diversidad reflejen la variedad de
fichas que el sistema encontrará en operación; sin él, cualquier métrica seguirá midiendo la
regla que generó esas etiquetas, como muestra la Sección~V-A. Solo una vez resueltos ambos resulta
informativa la evaluación con usuarios.

Una línea distinta, que este trabajo no aborda, es la detección del comportamiento fraudulento
del vendedor: determinar si un vendedor defrauda exige evidencia sobre operaciones consumadas
y una definición operativa de fraude que el comprador no puede establecer desde la ficha. Es
un problema de clasificación con verdad de referencia, no de estimación de riesgo previa a la
decisión, y abordarlo requeriría la colaboración de las plataformas o el acceso a denuncias
verificadas; constituye por ello una recomendación para investigaciones posteriores antes que
una extensión del sistema aquí descrito.
